% Introduction. Contract: ../references/section-playbook.md
% 420-650 words, four paragraphs, then 2-4 contribution items.
\section{Introduction}
\label{sec:introduction}

% P1 - problem
Certifying how a trained network will behave on unseen data is the point at
which a great deal of applied machine learning stops being measurable. Bounds
that are tight enough to be informative on realistic architectures have been
difficult to obtain directly~\cite{livni2020limitation, hellstrm2023generalization},
and the constructions that do succeed tend to route around the target model
rather than through it. Compression-based arguments certify a network by
certifying a smaller description of it~\cite{arora2018stronger, zhou2018nonvacuous,
shin2021compressing, bazinet2024sample}. Disagreement-based arguments go one
step further: they bound a tractable \surrogate, then bound how often the
\surrogate\ and the target disagree on unlabeled data, and combine the two into
a certificate for the target~\cite{jiang2021assessing, kirsch2022assessing}.
The appeal is that both halves are measurable even when the target is not.

% P2 - gap
The construction leaves one choice to the practitioner: which examples the
\surrogate\ is built from. Selection strategies for this role are well studied
in their own right, from uniform sampling through error-aware and
disagreement-aware criteria~\cite{coleman2019selection, xia2023refined,
benkert2024effective, feng2024tarot}, and they differ in how much control they
exert over the quantity the bound depends on. It would be natural to expect the
strategy with the most direct control over disagreement to produce the tightest
certificate. Reported comparisons do not consistently show that. However, the
evidence behind those comparisons mixes at least three sources of variation:
which data split was drawn, which selection method was used, and which
optimization path the training run happened to take. A single number cannot
separate them, so the observation is currently a result about a set of runs
rather than a result about a method.

% P3 - what this paper does
This paper measures the three factors apart. We freeze the \evalsem\ that give
the bound its meaning, register a comparison matrix before running it, and vary
the \splitseed\ and the \trajseed\ independently across the compared selection
methods. Contrasts are taken within a split, so the split effect cancels by
construction rather than by assumption, and a reporting threshold fixed before
execution decides whether a contrast is large enough to be called a finding at
all. Every run is preserved with a manifest identifying its method, seeds,
configuration, and exact command, so each reported statistic traces back to the
primary runs it aggregates.

% P4 - result and scope
Across 5 \splitseed{}s and 3 trajectory seeds per split and method, \methodrandom\
attains a lower mean \certbound\ than \methodmax\ by $0.011 \pm 0.005$. That
contrast does not clear the largest within-split trajectory spread in the
matrix, $0.032$, and its sign reverses at one of the five splits. The supported
conclusion is therefore that the two strategies are indistinguishable at this
scale, and that the trajectory, not the selection method, is the dominant
source of variation in the certificate. This is measured on one dataset at a
reduced experimental scale; it is not a claim about selection strategies in
general.

Our contributions are as follows:
\begin{itemize}[itemsep=0.5ex, parsep=0pt, topsep=0pt]
\item[\textbf{(1)}] A factorial design that separates \splitseed, selection
method, and \trajseed\ under frozen \evalsem, together with a reporting
threshold fixed in advance, which makes a method-level comparison identifiable
rather than confounded.

\item[\textbf{(2)}] A measurement of how large the trajectory-induced spread in
a disagreement-based certificate actually is, and the finding that it exceeds
the between-method contrast at this scale.

\item[\textbf{(3)}] A provenance-preserving pipeline in which every reported
aggregate enumerates the primary runs it was computed from, so any number in
the paper can be checked against a preserved record.
\end{itemize}
